\chapter{Thyroidectomy}

\section*{What Is This Surgery?}
A thyroidectomy is a surgical procedure involving the removal of all or part of the thyroid gland, an essential endocrine organ situated in the anterior neck, inferior to the larynx and anterior to the trachea. This surgery is primarily indicated for various thyroid disorders, including malignant tumors, suspicious nodules, symptomatic benign goiters, and hyperthyroidism that does not respond to medical therapy. The thyroid gland is crucial for metabolic regulation through the secretion of thyroid hormones, and its surgical excision requires careful consideration of potential impacts on hormonal balance and surrounding vital structures. The procedure aims to alleviate symptoms, prevent disease progression, or provide a definitive cure for malignancies, making it a significant intervention in endocrine surgery.

\section*{Alternative Names}
\begin{itemize}
  \item Total Thyroidectomy
  \item Hemithyroidectomy
  \item Thyroid Lobectomy
  \item Subtotal Thyroidectomy
  \item Near-Total Thyroidectomy
  \item Thyroid Gland Excision
\end{itemize}

\section*{Relevant Surgical Anatomy}
The thyroid gland is a butterfly-shaped organ composed of two lobes connected by an isthmus, located in the anterior neck. It is richly vascularized, receiving blood supply primarily from the superior thyroid artery (a branch of the external carotid artery) and the inferior thyroid artery (a branch of the thyrocervical trunk). Venous drainage occurs via the superior and middle thyroid veins into the internal jugular vein, and the inferior thyroid vein drains into the brachiocephalic vein.

The gland is closely associated with several critical structures, including the recurrent laryngeal nerves (RLNs), which innervate the vocal cords, and the parathyroid glands, which regulate calcium homeostasis. The RLNs typically run in the tracheoesophageal groove, making their identification crucial during surgery to prevent vocal cord paralysis. The parathyroid glands, usually four in number, are located on the posterior surface of the thyroid and must be preserved to avoid postoperative hypoparathyroidism.

Key anatomical landmarks include:
\begin{itemize}
  \item McBurney's point: Not applicable in thyroid surgery but important in other abdominal surgeries.
  \item Calot's triangle: Not directly related but represents the importance of understanding anatomical relationships in surgical dissection.
\end{itemize}

\section*{Surgical Principle \& Rationale}
The primary principle of thyroidectomy is the excision of diseased thyroid tissue while preserving surrounding critical structures, particularly the recurrent laryngeal nerves and parathyroid glands. The surgical goals include the complete removal of malignant or suspicious lesions with adequate margins, alleviation of compressive symptoms from large goiters, and resolution of hyperthyroidism by excising hyperactive gland tissue. 

The rationale behind this procedure is to achieve a definitive cure or effective control of thyroid pathologies. In cases of malignancy, the objective is to remove the primary tumor and any involved lymph nodes to prevent local recurrence and metastasis. For benign conditions, the surgery aims to decompress the trachea and esophagus, improving respiratory and swallowing functions. The extent of resection—whether total thyroidectomy or lobectomy—is determined by the specific indication, balancing the need for disease control with minimizing postoperative complications and the requirement for lifelong hormone replacement.

\section*{History \& Surgical Evolution}
The history of thyroid surgery dates back to ancient times, with early attempts often resulting in high mortality due to uncontrolled hemorrhage and infection. The 19th century marked significant advancements in the field. Theodor Billroth performed some of the first successful thyroidectomies in the 1870s, albeit with considerable mortality rates. The procedure was revolutionized by Emil Theodor Kocher, a Swiss surgeon, who standardized surgical techniques in the late 19th century. His emphasis on meticulous hemostasis, aseptic techniques, and careful anatomical dissection led to a dramatic reduction in mortality rates from over 40\% to less than 1\%, earning him the Nobel Prize in Physiology or Medicine in 1909.

The 20th century saw further advancements, including improved anesthetic techniques, a better understanding of parathyroid function, and the introduction of intraoperative nerve monitoring. The development of fine-needle aspiration biopsy (FNAB) in the latter half of the century allowed for more accurate preoperative diagnosis of thyroid nodules, guiding surgical decisions. Today, minimally invasive approaches and advanced imaging techniques continue to enhance the safety and efficacy of thyroidectomy.

\section*{Clinical Indications}
Thyroidectomy is indicated in the following clinical scenarios:

\begin{itemize}
  \item To definitively treat diagnosed thyroid cancer, including papillary, follicular, medullary, and anaplastic types.
  \item To remove thyroid nodules that are suspicious or indeterminate for malignancy based on fine-needle aspiration biopsy (Bethesda categories III-VI).
  \item To alleviate compressive symptoms such as dysphagia (difficulty swallowing), dyspnea (difficulty breathing), or hoarseness caused by a large benign goiter.
  \item To manage Graves' disease or toxic multinodular goiter that is refractory to medical therapy or radioactive iodine ablation.
  \item To remove a large substernal goiter that extends into the chest cavity, potentially causing airway obstruction or superior vena cava syndrome.
  \item To address cosmetic concerns related to a visibly enlarged thyroid gland in the neck.
  \item To prevent future complications in patients with certain genetic syndromes predisposing to medullary thyroid cancer (e.g., MEN 2).
\end{itemize}

\section*{Clinical Positioning}
Thyroidectomy is primarily a first-line treatment for most thyroid cancers and large, symptomatic benign goiters. In these cases, surgery is the initial intervention aimed at disease eradication or symptom resolution. For thyroid cancer, the extent of surgery (lobectomy vs. total thyroidectomy) is determined by tumor characteristics and staging. In cases of hyperthyroidism, thyroidectomy may be considered a secondary or tertiary treatment option, typically reserved for patients who have failed or cannot tolerate anti-thyroid medications or radioactive iodine therapy, or who have large goiters causing compressive symptoms in addition to hyperthyroidism. In advanced thyroid cancers with lymph node involvement, thyroidectomy may be part of a broader surgical approach that includes neck dissection, making it a component of a comprehensive primary treatment strategy.

\section*{Surgical Technique \& Procedure}
Thyroidectomy is performed under general anesthesia. The patient is positioned supine with the neck extended and a shoulder roll placed to optimize exposure of the anterior neck. A low, transverse cervical incision, often referred to as a Kocher incision, is made typically 1-2 cm above the sternal notch within a natural skin crease for optimal cosmetic outcome.

The surgical steps generally include:

\begin{itemize}
  \item **Incision and Dissection:** The skin, subcutaneous tissue, and platysma muscle are incised. Subplatysmal flaps are raised superiorly and inferiorly.
  
  \item **Strap Muscle Management:** The strap muscles (sternohyoid, sternothyroid) are identified and typically divided vertically in the midline or retracted laterally to expose the thyroid gland.
  
  \item **Middle Thyroid Vein Ligation:** The middle thyroid veins, which connect the thyroid gland to the internal jugular vein, are identified, ligated, and divided to mobilize the lateral aspect of the thyroid lobe.
  
  \item **Superior Pole Dissection:** The superior pole of the thyroid is carefully dissected. The superior thyroid artery and vein are ligated close to the gland to preserve the external branch of the superior laryngeal nerve, which controls vocal pitch.
  
  \item **Recurrent Laryngeal Nerve (RLN) Identification:** The RLN is meticulously identified and preserved. It typically runs in the tracheoesophageal groove, posterior to the thyroid gland. Intraoperative nerve monitoring may be used to confirm nerve function.
  
  \item **Parathyroid Gland Identification and Preservation:** The superior and inferior parathyroid glands, usually four in total, are identified and carefully dissected away from the thyroid gland, preserving their blood supply. If devascularized, they may be autotransplanted into a muscle pocket (e.g., sternocleidomastoid).
  
  \item **Inferior Pole Dissection:** The inferior thyroid artery and vein are ligated, again close to the gland, to preserve the blood supply to the parathyroid glands.
  
  \item **Gland Resection:** The thyroid lobe or entire gland is carefully separated from the trachea and removed.
  
  \item **Hemostasis and Closure:** Meticulous hemostasis is achieved. A small drain may be placed in the thyroid bed to prevent hematoma formation, though this is often omitted. The strap muscles are reapproximated, followed by closure of the platysma and skin layers.
\end{itemize}

\section*{Preoperative Workup \& Preparation}
Thorough preoperative preparation is essential to optimize patient safety and surgical outcomes. This typically involves:

\begin{itemize}
  \item **Medical Evaluation:** A comprehensive history and physical examination, including assessment of vocal cord mobility via indirect or direct laryngoscopy, especially if there is preoperative hoarseness or a large goiter.
  
  \item **Laboratory Tests:** Blood tests including complete blood count, electrolyte panel, renal and liver function tests, coagulation profile, and thyroid function tests (TSH, free T3, free T4). Serum calcium and parathyroid hormone (PTH) levels may be checked, particularly in cases of suspected parathyroid involvement or prior neck surgery.
  
  \item **Imaging Studies:** A neck ultrasound is standard for evaluating thyroid nodules and gland morphology. CT or MRI of the neck and chest may be performed for large or substernal goiters to assess tracheal deviation or compression.
  
  \item **Endocrine Optimization:** For patients with hyperthyroidism (e.g., Graves' disease, toxic multinodular goiter), it is crucial to achieve a euthyroid state preoperatively using anti-thyroid medications (e.g., methimazole, propylthiouracil). Lugol's iodine solution may be administered for 7-10 days prior to surgery to reduce gland vascularity and minimize intraoperative bleeding.
  
  \item **Anesthesia Clearance:** Patients undergo a preoperative anesthesia evaluation to assess fitness for general anesthesia and identify any potential risks.
  
  \item **Medication Review:** All medications are reviewed. Anticoagulants and antiplatelet agents may need to be held or bridged according to specific protocols.
  
  \item **NPO Status:** Patients are instructed to be NPO (nil per os) for a specified period before surgery, typically 6-8 hours, to prevent aspiration during anesthesia.
  
  \item **Informed Consent:** Detailed discussion of the procedure, potential risks, benefits, and alternatives, followed by obtaining informed consent.
  
  \item **Prophylactic Antibiotics:** A single dose of intravenous prophylactic antibiotics is typically administered prior to incision.
\end{itemize}

\section*{Contraindications \& Precautions}
**Absolute Contraindications:**

\begin{itemize}
  \item Uncorrectable severe coagulopathy or bleeding disorder.
  \item Severe, uncompensated medical comorbidities (e.g., acute myocardial infarction, uncontrolled sepsis) that render the patient unfit for general anesthesia and surgery.
  \item Patient refusal after thorough discussion of risks and benefits.
\end{itemize}

**Relative Contraindications and Precautions:**

\begin{itemize}
  \item Severe cardiovascular or pulmonary disease: Requires thorough preoperative optimization and careful anesthetic management.
  \item Active neck infection: Surgery should be delayed until the infection is resolved.
  \item Prior extensive neck radiation or surgery: Increases the risk of adhesions, nerve injury, and parathyroid damage due to altered anatomy and fibrosis.
  \item Uncontrolled hyperthyroidism: Increases the risk of thyroid storm; requires preoperative euthyroid state.
  \item Pregnancy: While not an absolute contraindication, surgery is typically deferred to the second trimester if possible, or medical management is preferred.
  \item Morbid obesity: May complicate surgical access and increase anesthetic risks.
\end{itemize}

\section*{Complications \& Risks}
While thyroidectomy is generally safe, potential risks and complications include:

\begin{itemize}
  \item **General Surgical Risks:**
    \begin{itemize}
      \item **Anesthesia Complications:** Allergic reactions, cardiac events, respiratory issues, nausea, vomiting.
      \item **Bleeding/Hematoma:** Accumulation of blood in the neck, which can cause swelling, pain, and potentially life-threatening airway compression requiring emergency re-exploration.
      \item **Infection:** Wound infection, though rare in clean neck surgery.
      \item **Pain:** Postoperative neck pain and discomfort.
    \end{itemize}
  
  \item **Procedure-Specific Risks:**
    \begin{itemize}
      \item **Recurrent Laryngeal Nerve (RLN) Injury:** Can be temporary or permanent. Unilateral injury causes hoarseness, dysphonia, and potential aspiration. Bilateral injury is rare but can lead to severe airway obstruction requiring tracheostomy.
      \item **Hypoparathyroidism:** Injury to or removal of the parathyroid glands can lead to transient or permanent hypocalcemia (low blood calcium). Symptoms include tingling (paresthesias) around the mouth and fingertips, muscle cramps, and in severe cases, tetany or seizures. Requires calcium and vitamin D supplementation.
      \item **Superior Laryngeal Nerve (SLN) Injury:** Specifically, injury to the external branch of the SLN can cause subtle changes in voice pitch, difficulty with high-pitched sounds, and vocal fatigue.
      \item **Thyroid Storm:** A rare but life-threatening complication in inadequately prepared hyperthyroid patients, characterized by fever, tachycardia, altered mental status, and multi-organ dysfunction.
      \item **Seroma:** Collection of serous fluid under the skin flap, usually resolves spontaneously or with aspiration.
      \item **Wound Complications:** Scarring (including keloid or hypertrophic scar formation), wound dehiscence (opening of the incision).
      \item **Chyle Leak:** Very rare injury to the thoracic duct (usually on the left side) resulting in lymphatic fluid leakage, typically managed conservatively but may require re-exploration.
      \item **Hypothyroidism:** After total thyroidectomy, lifelong thyroid hormone replacement (levothyroxine) is required. After lobectomy, some patients may still develop hypothyroidism.
    \end{itemize}
\end{itemize}

\section*{Postoperative Recovery Timeline}
Following thyroidectomy, patients are typically monitored in the Post-Anesthesia Care Unit (PACU) for the first few hours. Vital signs, airway patency, and neck swelling are continuously assessed. Pain management is initiated, and patients are monitored for signs of bleeding or airway compromise. Serum calcium levels are usually checked within 6-12 hours post-surgery, especially after total thyroidectomy, to detect early signs of hypocalcemia.

Over the first 24-48 hours, patients are encouraged to ambulate early. Diet is usually advanced from clear liquids to a regular diet as tolerated. Pain is managed with oral analgesics. Calcium and vitamin D supplementation may be initiated if hypocalcemia is detected or anticipated. Most patients are discharged within 1-2 days, provided they are stable, pain is controlled, and there are no signs of significant complications. Incision care instructions are provided, and patients are advised to avoid strenuous activities or heavy lifting for several weeks. Long-term, patients who undergo total thyroidectomy will require lifelong thyroid hormone replacement therapy.

\section*{Surgical Findings \& Pathology}
During thyroidectomy, the surgeon documents the extent of the resection, the condition of the thyroid tissue, and the status of surrounding structures, including the recurrent laryngeal nerves and parathyroid glands. The pathology report on the resected tissues provides definitive confirmation of the diagnosis, indicating whether the tissue is benign or malignant. For malignant cases, the report specifies the type of cancer (e.g., papillary, follicular, medullary), tumor size, presence of extrathyroidal extension, vascular invasion, surgical margin status, and involvement of any removed lymph nodes. These details are crucial for cancer staging and guiding subsequent management, including the need for radioactive iodine ablation or further surveillance.

For benign conditions, the pathology report confirms the absence of malignancy and characterizes the nature of the goiter or nodule. The discharge summary will outline the extent of surgery performed, any intraoperative findings (e.g., nerve integrity, parathyroid status), and initial postoperative course.

\section*{Expected Postoperative Course}
A normal, successful postoperative course following thyroidectomy is characterized by the resolution of preoperative symptoms, such as dysphagia and dyspnea, and the absence of complications. Patients typically experience a gradual decrease in pain and discomfort, with most reporting significant improvement within the first week. Serum calcium levels should stabilize within the normal range, indicating successful preservation of parathyroid function. Vocal cord function is assessed, and patients should not exhibit new or persistent hoarseness. The surgical incision should heal well without signs of infection, excessive swelling, or hematoma. For patients undergoing surgery for thyroid cancer, a normal result signifies successful resection of the tumor with clear margins, indicating a favorable prognosis and effective initial treatment.

\section*{Postoperative Care \& Follow-up}
Postoperative care is crucial for ensuring optimal recovery and monitoring for complications. Follow-up visits are typically scheduled 1-2 weeks after surgery for wound check, removal of any drains or sutures/staples, and initial discussion of the pathology report. 

Patients who undergo total thyroidectomy will require lifelong levothyroxine replacement therapy, initiated immediately post-op. For lobectomy patients, thyroid function is monitored, and replacement may be needed if hypothyroidism develops. 

Thyroid function tests, including serum TSH and free T4 levels, are checked 6-8 weeks post-surgery to adjust levothyroxine dosage and ensure appropriate thyroid hormone levels. Calcium and vitamin D levels are monitored until stable, and supplementation is adjusted or weaned as necessary.

For patients with malignancy, regular follow-up with endocrinologists is essential, including periodic neck ultrasounds, serum thyroglobulin levels, and anti-thyroglobulin antibodies to monitor for recurrence. Radioactive iodine ablation may be considered for high-risk cancers. 

Patients should be educated on warning signs to report, such as increasing neck swelling, difficulty breathing, fever, severe pain, or persistent tingling/numbness.

\section*{Epidemiology}
Thyroidectomy is one of the most common endocrine surgical procedures performed worldwide. The incidence of thyroid nodules is high, with up to 50\% of individuals having palpable nodules by age 60, though the vast majority are benign. Thyroid cancer, while relatively rare, has seen a rising incidence globally over the past few decades, largely attributed to increased detection of small, asymptomatic cancers due to widespread use of diagnostic imaging. It is more common in women than men, with a female-to-male ratio of approximately 3:1, and often affects individuals in their 30s to 50s. The most common indications for thyroidectomy include suspicious or malignant nodules, followed by symptomatic benign goiters and hyperthyroidism. While the mortality rate for thyroidectomy is exceedingly low (typically <0.1\%), the morbidity, particularly related to recurrent laryngeal nerve injury and hypoparathyroidism, remains a significant concern, with rates varying based on surgeon experience and extent of surgery.

\section*{Outcomes \& Prognosis}
The outcomes of thyroidectomy are generally excellent, with high success rates for resolving the underlying pathology. For benign conditions such as compressive goiter or hyperthyroidism, surgery typically provides definitive symptom resolution. For differentiated thyroid cancers (papillary and follicular), which constitute the majority of thyroid malignancies, the prognosis after appropriate surgical resection is highly favorable, with 10-year survival rates often exceeding 90-95\% for localized disease. Recurrence rates vary depending on the cancer type, stage, and initial treatment, but are generally low for low-risk cancers. Prognostic factors influencing outcomes include patient age, tumor size, histological subtype, presence of extrathyroidal extension, lymph node involvement, and distant metastases. While most patients achieve a cure, lifelong follow-up is essential for cancer surveillance and management of potential long-term complications such as hypothyroidism, which requires lifelong thyroid hormone replacement therapy.

\section*{References}
\begin{enumerate}
  \item MedlinePlus. Thyroidectomy. U.S. National Library of Medicine; 2024. Available from: \url{https://medlineplus.gov/ency/article/002931.htm}
  
  \item Brunicardi F, Andersen DK, Billiar TR, et al. Schwartz's Principles of Surgery. 11th ed. McGraw-Hill Education; 2019.
  
  \item American Thyroid Association (ATA) Guidelines Taskforce. Management of Thyroid Nodules and Differentiated Thyroid Cancer. Thyroid. 2015;25(1):1-133.
  
  \item P. L. G. Kocher. Ueber Kropfexstirpationen und ihre Folgen. Archiv f\"ur klinische Chirurgie. 1883;29:1-120.
\end{enumerate}

\clearpage